\documentclass[../DoAn.tex]{subfiles}
\begin{document}

\section{Đặt vấn đề}
\label{section:1.1}

Công tác thi đua, khen thưởng trong quân đội hiện được thực hiện theo Luật Thi đua, Khen thưởng số 06/2022/QH15 \cite{luat062022} cùng các nghị định và thông tư hướng dẫn.
Công việc này liên quan đến nhiều hồ sơ và lặp lại theo từng đợt xét, nên có thể ứng dụng công nghệ thông tin để giảm thao tác thủ công và sai sót. Tuy nhiên tại Học viện Khoa học Quân sự (HVKHQS), cán bộ vẫn quản lý thủ công bằng các văn bản, giấy tờ rời rạc chia sẻ qua mạng nội bộ.
Quá trình trao đổi và phê duyệt diễn ra trực tiếp, không có hệ thống quản lý tập trung, cũng không có kênh thông báo hay lưu vết tự động.
Cách làm đó bộc lộ những hạn chế chính, được khảo sát và phân tích kỹ hơn ở Mục~\ref{section:2.1}.

Thứ nhất, các danh hiệu cá nhân hằng năm xếp thành nhiều cấp nối tiếp nhau, trong đó danh hiệu cấp trên lấy danh hiệu cấp dưới làm điều kiện.
Ở dưới cùng và được xét hằng năm là Chiến sĩ thi đua cơ sở (CSTĐCS).
Để đạt các danh hiệu cao hơn như Bằng khen của Bộ trưởng Bộ Quốc phòng (BKBQP), Chiến sĩ thi đua toàn quân (CSTĐTQ) hay Bằng khen của Thủ tướng Chính phủ (BKTTCP), quân nhân vừa phải giữ CSTĐCS liên tục đủ số năm quy định, vừa phải tích lũy đủ số lần các danh hiệu ở cấp dưới và có thành tích nghiên cứu khoa học hằng năm.
Việc xét càng phức tạp khi chỉ những danh hiệu đạt được trong ít năm gần nhất mới được tính, còn các danh hiệu cũ hơn không còn giá trị trong lần xét hiện tại.
Vì vậy, để xác định một quân nhân đã đủ điều kiện hay chưa, cán bộ phải tra cứu và đối chiếu kỹ lịch sử khen thưởng của từng người, tốn nhiều thời gian và tiềm ẩn nguy cơ bỏ sót.

Thứ hai, công tác khen thưởng ở Học viện gồm bảy nhóm nghiệp vụ, mỗi nhóm có tiêu chuẩn xét riêng.
Bên cạnh khen thưởng cá nhân và đơn vị hằng năm, Học viện còn xét Huy chương Chiến sĩ vẻ vang (HCCSVV) theo mốc 10, 15, 20 năm phục vụ; Huân chương Bảo vệ Tổ quốc (HCBVTQ) theo thời gian giữ chức vụ; Kỷ niệm chương, hiện có Kỷ niệm chương vì sự nghiệp xây dựng Quân đội nhân dân Việt Nam (KNC), theo số năm công tác với mốc khác nhau giữa nam và nữ; và Huy chương Quân kỳ quyết thắng (HCQKQT) cho người phục vụ đủ 25 năm. Cuối cùng, thành tích nghiên cứu khoa học (NCKH) được quản lý thành một nhóm riêng, đồng thời là điều kiện bắt buộc khi xét các danh hiệu cấp cao trong khen thưởng cá nhân hằng năm.

Dữ liệu của các nhóm này lại nằm rải rác trong nhiều tệp khác nhau, nên khó kiểm tra và khó giữ được sự nhất quán.

Thứ ba, các tệp dùng chung không có cơ chế lưu vết thao tác (audit log).
Khi một bản ghi bị sửa hay xoá, hệ thống không lưu lại ai thực hiện, vào thời điểm nào và vì lý do gì, khiến việc truy vết, rà soát về sau gặp nhiều khó khăn.
Việc tiếp cận dữ liệu phụ thuộc vào người hay đơn vị đang giữ tệp, không có cơ chế phân quyền theo cơ cấu tổ chức thực tế của Học viện gồm cấp Phòng, cấp Khoa và từng quân nhân.

Những hạn chế trên, từ việc xét các điều kiện chuỗi phức tạp, dữ liệu bảy nhóm nghiệp vụ nằm rải rác, đến việc thiếu cơ chế lưu vết và phân quyền sát với cơ cấu tổ chức, cho thấy cách quản lý thủ công hiện nay không còn phù hợp với công tác khen thưởng tại Học viện. Đây chính là bài toán mà đồ án hướng tới giải quyết.

\section{Mục tiêu và phạm vi đề tài}
\label{section:1.2}

Trên thị trường có các phần mềm quản lý nhân sự nói chung, nhưng chưa có sản phẩm phù hợp với đặc thù nghiệp vụ khen thưởng theo chuỗi điều kiện nhiều năm và cơ cấu tổ chức của Học viện.
Để khắc phục những hạn chế nêu trên, mục tiêu chính của đồ án là xây dựng một hệ thống web hỗ trợ Phòng Chính trị Học viện Khoa học Quân sự thực hiện các nhóm nhiệm vụ sau:

\begin{enumerate}
    \item \textbf{Quản lý hồ sơ và danh mục cơ bản}: quân nhân, đơn vị theo cơ cấu hai cấp gồm Cơ quan đơn vị và Đơn vị trực thuộc, chức vụ và lịch sử biến động chức vụ.
    \item \textbf{Đề xuất và phê duyệt khen thưởng}: tạo đợt đề xuất cho bảy nhóm danh hiệu với bộ lọc hỗ trợ tra cứu nhanh theo đơn vị, năm, giới tính và mốc thời gian phục vụ; quy trình phê duyệt nhiều bước, lưu tệp PDF quyết định kèm số quyết định. Riêng khen thưởng đột xuất phát sinh ngoài chu kỳ được ghi nhận trực tiếp dưới dạng quyết định, không qua bước đề xuất.
    \item \textbf{Tự động hoá nghiệp vụ}: kiểm tra điều kiện chuỗi danh hiệu cho cá nhân (BKBQP, CSTĐTQ, BKTTCP) và đơn vị (BKBQP, BKTTCP), tự động tạo thông báo gợi ý đủ hoặc chưa đủ điều kiện kèm lý do bằng tiếng Việt, nhập lịch sử khen thưởng từ Excel theo cơ chế xem trước rồi xác nhận, và xuất danh sách theo nhiều tiêu chí. Quá trình nhập tuân thủ tính toàn vẹn của giao dịch cơ sở dữ liệu: hoặc ghi nhận toàn bộ, hoặc hoàn tác hoàn toàn.
    \item \textbf{Bảo mật và truy vết}: phân quyền theo bốn cấp tương ứng đặc thù tổ chức. Quản trị viên cấp cao quản trị hạ tầng, Cán bộ Phòng Chính trị toàn quyền nghiệp vụ, Chỉ huy đơn vị đề xuất trong phạm vi đơn vị, còn Người dùng tra cứu hồ sơ cá nhân. Toàn bộ thao tác làm thay đổi dữ liệu đều được ghi vào nhật ký hệ thống.
    \item \textbf{Vận hành}: sao lưu cơ sở dữ liệu định kỳ và quản lý nhật ký hệ thống.
\end{enumerate}

\textbf{Phạm vi đề tài} giới hạn trong các nhiệm vụ trực tiếp phục vụ công tác xét khen thưởng.
Phần mềm vận hành tập trung trên máy chủ của Học viện, truy cập qua mạng nội bộ và không công khai ra Internet công cộng, phù hợp với đặc thù bảo mật của đơn vị quân đội.
Đồ án không bao gồm các nghiệp vụ tiền lương, phụ cấp, hồ sơ y tế, đánh giá cán bộ định kỳ và các quy trình hành chính khác.
Việc tích hợp với hệ thống quản lý nhân sự tổng thể của Bộ Quốc phòng cũng nằm ngoài phạm vi do vượt khỏi trọng tâm khen thưởng.

\section{Định hướng giải pháp}
\label{section:1.3}

Phần mềm được xây dựng theo mô hình máy khách/máy chủ, gồm hai thành phần triển khai độc lập, có thể đặt chung trên một máy chủ hoặc tách ra nhiều máy.
Phần giao diện dùng Next.js theo App Router viết bằng TypeScript, trong đó Ant Design đảm nhiệm lớp giao diện chính và Tailwind CSS bổ sung cho bố cục.
Phần xử lý dùng Express trên Node.js, cũng viết bằng TypeScript và tổ chức theo kiến trúc phân tầng tách biệt định tuyến, xử lý nghiệp vụ và truy cập dữ liệu.
Dữ liệu được lưu trong PostgreSQL và truy cập qua Prisma ORM với lược đồ khai báo tập trung, áp dụng cơ chế cập nhật lược đồ tự động (Prisma Migration).
Hệ thống xác thực bằng cặp JSON Web Token có cơ chế làm mới luân phiên, đồng thời kiểm tra vai trò người dùng trước mỗi thao tác cần quyền truy cập.
Các sự kiện phê duyệt, từ chối, xoá đề xuất và nhập Excel khối lượng lớn được thông báo theo thời gian thực qua Socket.IO.

Đóng góp chính của đồ án nằm ở việc tự động hoá nghiệp vụ xét khen thưởng, vốn có nhiều điều kiện phụ thuộc theo thời gian và dễ phát sinh sai lệch khi xử lý thủ công.
Nổi bật là cơ chế xét điều kiện chuỗi danh hiệu thống nhất cho cả cá nhân và đơn vị. Bên cạnh đó, bảy nhóm khen thưởng vốn rời rạc được hợp nhất vào một quy trình chung từ đề xuất, phê duyệt đến gắn quyết định. Cuối cùng, nhật ký hệ thống ghi vết mọi thao tác làm thay đổi dữ liệu.
Các đóng góp này được phân tích chi tiết trong Chương~\ref{chapter:SolutionAndContribution}.

\section{Bố cục đồ án}
\label{section:1.4}

Báo cáo gồm sáu chương. Ngoài chương này, phần còn lại được tổ chức như sau.
Chương~\ref{chapter:Related_works} trình bày khảo sát hiện trạng, biểu đồ use case, đặc tả các use case trọng yếu và các yêu cầu phi chức năng của hệ thống.
Chương~\ref{chapter:Methodology} giới thiệu các thành phần công nghệ được lựa chọn cùng lý do lựa chọn trong ngữ cảnh phần mềm quản lý khen thưởng.
Chương~\ref{chapter:Experiment} trình bày kiến trúc phân tầng, thiết kế cơ sở dữ liệu, các biểu đồ gói, lớp và tuần tự, kết quả kiểm thử cùng hướng dẫn triển khai.
Chương~\ref{chapter:SolutionAndContribution} phân tích năm đóng góp nổi bật, mỗi đóng góp gồm thực trạng, giải pháp và kết quả đạt được.
Cuối cùng, Chương~\ref{chapter:conclusion} tổng kết những kết quả đã đạt được và đề xuất các hướng phát triển tiếp theo.

Chương này đã trình bày bối cảnh và động lực của đề tài. Xuất phát từ thực trạng quản lý khen thưởng thủ công tại Học viện Khoa học Quân sự cùng những hạn chế của cách làm hiện tại, đồ án xác định mục tiêu, phạm vi và định hướng giải pháp công nghệ cho phần mềm.
Những nội dung này là cơ sở để chương tiếp theo đi sâu khảo sát hiện trạng, mô hình hoá yêu cầu qua biểu đồ use case và xác định các yêu cầu phi chức năng của hệ thống.

\end{document}
